\documentclass{subfiles}
\begin{document}
    Let us consider the following: assume we are given a data structure,
        and we are asked to compute a sequence of \(m\) operations on it;
        what is the overall complexity of all \(m\) operations, i.e,
        what is the time needed to compute all \(m\) operations.

    It follows easily that, applying a worst case analysis,
        we need \(m \cdot T_{MAX}(n)\), where \(T_{MAX}(n)\) 
        is the time complexity of the most expesive operation.
        For instance, under the worst case assumption,
        a sequence of \(m\) deletion from a tree has a complexity of \(\bigO{m \cdot n}\).

    We now introduce the notion of \emph{amortized analysis},
        and show how, over a sequence of operations, 
        the real time complexity of an algorithm differs from the worst case one.

    In its essence, amortized analysis does nothing more than averaging the complexity of 
        all operations. That is, if \(T_{1}(n)\) is the complexity of the first operation,
        \(T_{2}(n)\) the one for the second operation, \(\ldots\), 
        \(T_{m}(n)\) the complexity of the \(m\)-th operation, 
        then the amortized time is given as 
        \[
            T_{amortized}(n) = \frac{1}{n} \sum_{i = 1}^{m} T_{i}(n).
        \]
        In other words, amortized time represents the time needed for each of 
        the \(m\) operations.

    The next section describes a simple ways to perform amortized analysis.
    \subsubsection{Accounting method}
    \subfile{../subsub/1.2.1 - Accounting method}
\end{document}
